%----------------------------------------------------------------------------------------
% INVERSE, DETERMINANT EN RANG. Alle opgaven zijn eigen werk.
%
% Opbouw: Start, leerstof (theorie, werkwijze of uitleg, voorbeeld, korte oefening), en op een
% nieuwe pagina de Opdrachten (door elkaar, stijgende moeilijkheid), Koppel en Test jezelf.
% Labels van dit hoofdstuk beginnen met h3-.
%----------------------------------------------------------------------------------------

\mychapter{3}{Inverse, determinant en rang}

\lkr{Leerplan B+S'': LPD 36 (de rang van matrices, de inverse van inverteerbare matrices en de determinant van vierkante matrices), met de wenken: de inverse via $[M \,|\, I] \to [I \,|\, M^{-1}]$, de determinant ingeleid met de inverse van een $2 \times 2$-matrix, ontwikkelen naar een rij of kolom en de regel van Sarrus, en het verband tussen rang, inverteerbaarheid, determinant en stelsels. De regel van Cramer (wenk bij LPD 37) staat bij de determinant. Grote berekeningen mogen met het rekentoestel.}

Een vergelijking $3x = 12$ los je op door beide leden met $\frac13$ te vermenigvuldigen, het omgekeerde van 3. Kan dat ook met een matrixvergelijking $AX = B$? Delen door een matrix bestaat niet, maar soms is er een matrix die de rol van $\frac13$ speelt: de inverse. In dit hoofdstuk leer je wanneer die bestaat en hoe je ze berekent. Een enkel getal, de determinant, zegt meteen of ze bestaat. En de rang telt hoeveel echt verschillende vergelijkingen er in een stelsel zitten.

\begin{oefeningen}[Start]
\oef{h3-start-1}

\oef{h3-verken}
\end{oefeningen}

\lkr{Laat oefening \ref{oef:h3-verken} maken vóór de uitleg. In a vinden leerlingen dat $AB = BA = I$. In b lossen ze een stelsel op door met $B$ te vermenigvuldigen, zonder rijherleiden. In c merken ze dat het voor sommige matrices niet lukt: de tweede rij is het dubbel van de eerste.}

%========================================================================================
\section{De inverse matrix}

\subsection{Definitie}

Bij getallen is $\frac13$ het omgekeerde van 3, omdat $3 \cdot \frac13 = 1$. Bij matrices speelt de eenheidsmatrix $I$ de rol van 1. Omdat het product van matrices niet commutatief is, eisen we dat het in beide volgordes lukt.

\begin{theorie}
Een vierkante matrix $A$ is \textbf{\textcolor{accent}{inverteerbaar}} als er een matrix $B$ bestaat met \quad $\formule{AB = BA = I}$

\vspace{1mm}
Die matrix heet de \textbf{\textcolor{accent}{inverse}} van $A$, en je noteert ze $\formule{A^{-1}}$.

\begin{stelling}{De inverse is uniek}
Een vierkante matrix heeft hoogstens één inverse.
\end{stelling}

\begin{bewijs}
Stel dat $B$ en $C$ allebei inversen zijn van $A$. Met de associativiteit van het product is
\bewijsregel{B = BI = B(AC) = (BA)C = IC = C .}
\end{bewijs}
\end{theorie}

Pas door die stelling mag je spreken over \textit{de} inverse, en heeft de notatie $A^{-1}$ zin. Niet elke vierkante matrix is inverteerbaar, net zoals 0 geen omgekeerde heeft. In de Verken lukte het niet voor $\km{1 & 2 \\ 2 & 4}$. En in hoofdstuk 1 zag je matrices $A \neq O$ en $B \neq O$ met $AB = O$: zo'n $A$ heeft geen inverse, want anders was $B = A^{-1}AB = A^{-1}O = O$.

\begin{oefeningen}
\oef{h3-isinverse}
\end{oefeningen}

\subsection{Matrixvergelijkingen}

Is $A$ inverteerbaar, dan los je $AX = B$ op zoals $3x = 12$: vermenigvuldig beide leden met $A^{-1}$. Omdat de volgorde telt, vermenigvuldig je aan dezelfde kant als waar $A$ staat.

\begin{theorie}
\begin{stelling}{Matrixvergelijking oplossen}
Voor een inverteerbare matrix $A$ is

\vspace{1mm}
\qquad $\formule{AX = B \iff X = A^{-1}B}$ \qquad en \qquad $\formule{XA = B \iff X = BA^{-1}}$
\end{stelling}

\begin{bewijs}
Stel $AX = B$. Vermenigvuldig beide leden links met $A^{-1}$:
\bewijsregel{A^{-1}AX = A^{-1}B \quad\Rightarrow\quad IX = A^{-1}B \quad\Rightarrow\quad X = A^{-1}B .}
Omgekeerd is $X = A^{-1}B$ een oplossing, want $A(A^{-1}B) = (AA^{-1})B = IB = B$. Voor $XA = B$ vermenigvuldig je rechts met $A^{-1}$: $XAA^{-1} = BA^{-1}$, dus $X = BA^{-1}$, en omgekeerd is $(BA^{-1})A = B(A^{-1}A) = B$.
\end{bewijs}
\end{theorie}

Er is dus precies één oplossing. De inverse komt aan de kant waar $A$ stond: links bij $AX = B$, rechts bij $XA = B$. Kies je de verkeerde kant, dan krijg je een andere matrix, of een product dat niet eens bestaat: bij een kolom $B$ en een $2 \times 2$-matrix $A$ is $BA^{-1}$ een $2 \times 1$ maal een $2 \times 2$.

\voorbeeld{}{Los op: $\begin{cases} 2x + y = 4 \\ 5x + 3y = 11 \end{cases}$, wetende dat $\km{2 & 1 \\ 5 & 3}^{-1} = \km{3 & -1 \\ -5 & 2}$.\label{vb:h3-vgl}}{%
\uitleg{Het stelsel is $AX = B$. Omdat we $A^{-1}$ kennen, is de oplossing gewoon een product.}
\begin{stappen}
X = A^{-1}B = \begin{bmatrix} 3 & -1 \\ -5 & 2 \end{bmatrix} \begin{bmatrix} 4 \\ 11 \end{bmatrix} & Links vermenigvuldigen met $A^{-1}$. \\[3.5ex]
X = \begin{bmatrix} 12 - 11 \\ -20 + 22 \end{bmatrix} = \begin{bmatrix} 1 \\ 2 \end{bmatrix} & Dus $x = 1$ en $y = 2$. Controle: $2 + 2 = 4$ en $5 + 6 = 11$.
\end{stappen}}

Dat is vooral handig als je hetzelfde stelsel met veel verschillende rechterleden moet oplossen: je berekent $A^{-1}$ één keer.

\begin{oefeningen}
\oef{h3-vglkort}
\end{oefeningen}

Staat de onbekende matrix links van $A$, dan vermenigvuldig je rechts met $A^{-1}$.

\voorbeeld{}{Bepaal $X$ als $X \begin{bmatrix} 2 & 1 \\ 5 & 3 \end{bmatrix} = \begin{bmatrix} 1 & 1 \\ 0 & 2 \end{bmatrix}$. Gebruik de inverse uit het vorige voorbeeld.\label{vb:h3-xa}}{%
\uitleg{De vergelijking heeft de vorm $XA = B$: $A$ staat rechts van $X$. Dus komt $A^{-1}$ rechts van $B$.}
\begin{stappen}
X = BA^{-1} = \begin{bmatrix} 1 & 1 \\ 0 & 2 \end{bmatrix} \begin{bmatrix} 3 & -1 \\ -5 & 2 \end{bmatrix} & Rechts vermenigvuldigen met $A^{-1}$. \\[3.5ex]
X = \begin{bmatrix} 3 - 5 & -1 + 2 \\ -10 & 4 \end{bmatrix} = \begin{bmatrix} -2 & 1 \\ -10 & 4 \end{bmatrix} & Rij maal kolom. \\[3.5ex]
XA = \begin{bmatrix} -4 + 5 & -2 + 3 \\ -20 + 20 & -10 + 12 \end{bmatrix} = \begin{bmatrix} 1 & 1 \\ 0 & 2 \end{bmatrix} & Controle: dat is $B$.
\end{stappen}
Wie $A^{-1}B$ berekent, vindt $\km{3 & 1 \\ -5 & -1}$, en dat is geen oplossing.}

\begin{werkwijze}
Kijk aan welke kant van $X$ de matrix $A$ staat. Vermenigvuldig beide leden aan die kant met $A^{-1}$. Controleer door $X$ in te vullen in de oorspronkelijke vergelijking.
\end{werkwijze}

\begin{oefeningen}
\oef{h3-matrixvgl}
\end{oefeningen}

\subsection{De inverse berekenen}

Hoe vind je $A^{-1}$? Bij een $3 \times 3$-matrix zijn de kolommen van $A^{-1}$ de oplossingen van drie stelsels: $AX = $ de eerste, tweede en derde kolom van $I$. Die stelsels hebben dezelfde coëfficiëntenmatrix, dus rijherleid je ze in één keer.

\begin{theorie}
\begin{stelling}{Inverse met Gauss-Jordan}
Een vierkante matrix $A$ is inverteerbaar als en slechts als $A$ rijherleidt tot $I$. Dan rijherleidt \quad $\formule{[A \,|\, I] \ \to \ [I \,|\, A^{-1}]}$
\end{stelling}
\end{theorie}

Waarom staat rechts $A^{-1}$? Rijherleiden van $[A \,|\, I]$ lost tegelijk de stelsels $AX = $ kolom $j$ van $I$ op. Staat links $I$, dan staat rechts in kolom $j$ de oplossing van dat stelsel. Noem $C$ de matrix rechts; dan is kolom $j$ van $AC$ gelijk aan kolom $j$ van $I$, dus $AC = I$. Dat ook $CA = I$, en dat $A$ niet inverteerbaar is als je links geen $I$ krijgt, bewijzen we niet: dat vraagt meer theorie dan deze cursus biedt. Krijg je links een nulrij, dan stop je dus: $A$ is niet inverteerbaar.

\voorbeeld{}{Bereken de inverse van $A = \begin{bmatrix} 1 & 2 & 1 \\ 0 & 1 & 1 \\ 1 & 2 & 2 \end{bmatrix}$.\label{vb:h3-inverse}}{%
\uitleg{We zetten $I$ naast $A$ en rijherleiden tot links $I$ staat. Elke rijbewerking doen we op de hele rij, ook rechts van de streep.}
\begin{stappenbreed}
\um{ccc|ccc}{1 & 2 & 1 & 1 & 0 & 0 \\ 0 & 1 & 1 & 0 & 1 & 0 \\ 1 & 2 & 2 & 0 & 0 & 1} & De matrix $[A \,|\, I]$. \\[5ex]
\um{ccc|ccc}{1 & 0 & -1 & 1 & -2 & 0 \\ 0 & 1 & 1 & 0 & 1 & 0 \\ 0 & 0 & 1 & -1 & 0 & 1} \ \begin{array}{l} \rb{R_1 - 2R_2} \\ \\ \rb{R_3 - R_1} \end{array} & $R_3 - R_1$ met de oorspronkelijke $R_1$. \\[5ex]
\um{ccc|ccc}{1 & 0 & 0 & 0 & -2 & 1 \\ 0 & 1 & 0 & 1 & 1 & -1 \\ 0 & 0 & 1 & -1 & 0 & 1} \ \begin{array}{l} \rb{R_1 + R_3} \\ \rb{R_2 - R_3} \\ \\ \end{array} & Links staat $I$. \\[5ex]
A^{-1} = \begin{bmatrix} 0 & -2 & 1 \\ 1 & 1 & -1 \\ -1 & 0 & 1 \end{bmatrix} & Controle: de eerste rij van $A$ maal de eerste kolom van $A^{-1}$ is $0 + 2 - 1 = 1$, enzovoort.
\end{stappenbreed}}

\begin{werkwijze}
Schrijf $[A \,|\, I]$ en rijherleid zoals bij Gauss-Jordan. Staat links $I$, dan staat rechts $A^{-1}$. Krijg je links een nulrij, dan is $A$ niet inverteerbaar. Controleer met $AA^{-1} = I$, minstens voor één rij en kolom.
\end{werkwijze}

\nujij{Bereken met de inverse uit voorbeeld \ref{vb:h3-inverse} de oplossing van $AX = \km{4 \\ 1 \\ 5}$.}{$X = A^{-1}\km{4 \\ 1 \\ 5} = \km{0 - 2 + 5 \\ 4 + 1 - 5 \\ -4 + 0 + 5} = \km{3 \\ 0 \\ 1}$.}

\begin{oefeningen}
\oef{h3-inverse}
\end{oefeningen}

\subsection{Rekenregels voor de inverse}

Bij getallen is het omgekeerde van $\frac13$ opnieuw 3, en is $\frac{1}{ab} = \frac1a \cdot \frac1b$. Voor matrices geldt de eerste regel ook, maar bij de tweede speelt de volgorde een rol.

\begin{theorie}
\begin{stelling}{Inverse van de inverse}
Is $A$ inverteerbaar, dan is ook $A^{-1}$ inverteerbaar, en \qquad $\formule{(A^{-1})^{-1} = A}$
\end{stelling}

\begin{bewijs}
Uit de definitie van $A^{-1}$ volgt $A^{-1}A = AA^{-1} = I$. Lees je dat met $A^{-1}$ als de matrix en $A$ als de kandidaat-inverse, dan staat er net dat $A$ een inverse is van $A^{-1}$. Omdat de inverse uniek is, is $(A^{-1})^{-1} = A$.
\end{bewijs}
\end{theorie}

Een matrix en haar inverse zijn dus elkaars inverse: wie twee keer omkeert, komt terug bij het begin.

\begin{theorie}
\begin{stelling}{Inverse van een product}
Zijn $A$ en $B$ inverteerbare $n \times n$-matrices, dan is ook $AB$ inverteerbaar, en

\vspace{1mm}
\qquad $\formule{(AB)^{-1} = B^{-1}A^{-1}}$
\end{stelling}

\begin{bewijs}
We gaan na dat $B^{-1}A^{-1}$ aan de definitie voldoet. Met de associativiteit is
\bewijsregel{(AB)(B^{-1}A^{-1}) = A(BB^{-1})A^{-1} = AIA^{-1} = AA^{-1} = I}
en
\bewijsregel{(B^{-1}A^{-1})(AB) = B^{-1}(A^{-1}A)B = B^{-1}IB = B^{-1}B = I .}
\end{bewijs}
\end{theorie}

Let op de volgorde: zoals bij kousen en schoenen. Wie eerst kousen en dan schoenen aandoet, trekt eerst de schoenen uit. Omdat het product niet commutatief is, is $A^{-1}B^{-1}$ meestal niet de inverse van $AB$.

\voorbeeld{}{Bereken $(AB)^{-1}$ met de inversen van $A = \begin{bmatrix} 2 & 1 \\ 1 & 1 \end{bmatrix}$ en $B = \begin{bmatrix} 1 & 2 \\ 0 & 1 \end{bmatrix}$.\label{vb:h3-invproduct}}{%
\uitleg{We hoeven $AB$ niet eerst uit te rekenen: de inverse is het product van de inversen, in omgekeerde volgorde.}
\begin{stappen}
A^{-1} = \begin{bmatrix} 1 & -1 \\ -1 & 2 \end{bmatrix}, \quad B^{-1} = \begin{bmatrix} 1 & -2 \\ 0 & 1 \end{bmatrix} & Beide determinanten zijn 1: verwissel de diagonaal en verander het teken van de rest. \\[3.5ex]
(AB)^{-1} = B^{-1}A^{-1} = \begin{bmatrix} 1 & -2 \\ 0 & 1 \end{bmatrix} \begin{bmatrix} 1 & -1 \\ -1 & 2 \end{bmatrix} & Eerst $B^{-1}$, dan $A^{-1}$. \\[3.5ex]
= \begin{bmatrix} 1 + 2 & -1 - 4 \\ -1 & 2 \end{bmatrix} = \begin{bmatrix} 3 & -5 \\ -1 & 2 \end{bmatrix} & Rij maal kolom. \\[3.5ex]
AB = \begin{bmatrix} 2 & 5 \\ 1 & 3 \end{bmatrix}, \quad AB \cdot \begin{bmatrix} 3 & -5 \\ -1 & 2 \end{bmatrix} = \begin{bmatrix} 1 & 0 \\ 0 & 1 \end{bmatrix} & Controle.
\end{stappen}
In de verkeerde volgorde is $A^{-1}B^{-1} = \km{1 & -3 \\ -1 & 4}$: dat is niet de inverse.}

\begin{oefeningen}
\oef{h3-k-invregels}
\end{oefeningen}

%========================================================================================
\section{De determinant}

\subsection{De determinant van een \texorpdfstring{$2 \times 2$}{2x2}-matrix}

Zoek je met $[A \,|\, I]$ de inverse van een algemene matrix $A = \km{a & b \\ c & d}$, dan komt telkens hetzelfde getal $ad - bc$ in de noemer. Dat getal beslist of de inverse bestaat.

\begin{theorie}
De \textbf{\textcolor{accent}{determinant}} van $A = \begin{bmatrix} a & b \\ c & d \end{bmatrix}$ is \quad $\formule{\det A = \begin{vmatrix} a & b \\ c & d \end{vmatrix} = ad - bc}$

\begin{stelling}{Inverse van een $2 \times 2$-matrix}
Is $\det A \neq 0$, dan is \quad $\formule{A^{-1} = \dfrac{1}{ad - bc} \begin{bmatrix} d & -b \\ -c & a \end{bmatrix}}$
\end{stelling}

\begin{bewijs}
Reken de twee producten uit:
\bewijsregel{\begin{bmatrix} a & b \\ c & d \end{bmatrix} \begin{bmatrix} d & -b \\ -c & a \end{bmatrix} = \begin{bmatrix} ad - bc & -ab + ba \\ cd - dc & -cb + da \end{bmatrix} = (ad - bc) I}
\bewijsregel{\begin{bmatrix} d & -b \\ -c & a \end{bmatrix} \begin{bmatrix} a & b \\ c & d \end{bmatrix} = \begin{bmatrix} da - bc & db - bd \\ -ca + ac & -cb + ad \end{bmatrix} = (ad - bc) I}
Na vermenigvuldigen met $\frac{1}{ad - bc}$ staat er telkens $I$.
\end{bewijs}
\end{theorie}

De inverse van een $2 \times 2$-matrix schrijf je dus zo op: verwissel de getallen op de hoofddiagonaal, verander het teken van de andere twee, en deel door de determinant. Is $\det A = 0$, dan is $A$ niet inverteerbaar: dat bewijs je in opdracht \ref{oef:h3-o11}.

\voorbeeld{}{Bepaal de inverse van $A = \begin{bmatrix} 3 & 1 \\ 2 & -1 \end{bmatrix}$.\label{vb:h3-inverse2}}{%
\uitleg{We berekenen eerst de determinant. Is die niet 0, dan gebruiken we de formule.}
\begin{stappen}
\det A = 3 \cdot (-1) - 1 \cdot 2 = -5 & Niet 0: $A$ is inverteerbaar. \\[1ex]
A^{-1} = \frac{1}{-5}\begin{bmatrix} -1 & -1 \\ -2 & 3 \end{bmatrix} = \begin{bmatrix} \frac15 & \frac15 \\[.5ex] \frac25 & -\frac35 \end{bmatrix} & Diagonaal verwisselen, de andere twee van teken veranderen, delen door $-5$.
\end{stappen}}

\begin{oefeningen}
\oef{h3-inverse2k}
\end{oefeningen}

\subsection{De determinant van een \texorpdfstring{$3 \times 3$}{3x3}-matrix}

Ook bij grotere vierkante matrices bestaat een determinant, een getal dat zegt of de matrix inverteerbaar is. Je berekent ze door ze terug te brengen tot determinanten van $2 \times 2$-matrices.

\begin{theorie}
De \textbf{\textcolor{accent}{minor}} $M_{ij}$ van $A$ is de determinant van de matrix die overblijft als je rij $i$ en kolom $j$ schrapt. De determinant van een $3 \times 3$-matrix $A$, \textbf{\textcolor{accent}{ontwikkeld naar de eerste rij}}, is

\vspace{1mm}
\qquad $\formule{\det A = a_{11}M_{11} - a_{12}M_{12} + a_{13}M_{13}}$

\begin{stelling}{Ontwikkelen naar een rij of kolom}
Je mag ook naar een andere rij of kolom ontwikkelen, met de tekens volgens het schaakbordpatroon $\begin{smallmatrix} + & - & + \\ - & + & - \\ + & - & + \end{smallmatrix}$: de uitkomst is altijd dezelfde.
\end{stelling}
\end{theorie}

Die stelling bewijzen we niet in het algemeen. Met de zes termen van de regel van Sarrus hieronder kun je ze wel narekenen: elke rij en elke kolom geeft dezelfde zes termen. Kies dus een rij of kolom met veel nullen: elke 0 is een term minder.

\voorbeeld{}{Bereken $\det B$ voor $B = \begin{bmatrix} 2 & -1 & 3 \\ 1 & 0 & 2 \\ 4 & 1 & -1 \end{bmatrix}$.\label{vb:h3-det}}{%
\uitleg{De tweede rij bevat een 0, dus ontwikkelen we naar die rij. De tekens in rij 2 zijn $-$, $+$, $-$.}
\begin{stappenbreed}
\det B = -1 \cdot \begin{vmatrix} -1 & 3 \\ 1 & -1 \end{vmatrix} + 0 - 2 \cdot \begin{vmatrix} 2 & -1 \\ 4 & 1 \end{vmatrix} & Schrap telkens de rij en de kolom van het element. \\[4ex]
= -1 \cdot (1 - 3) - 2 \cdot (2 + 4) = 2 - 12 = -10 & Twee determinanten van $2 \times 2$-matrices.
\end{stappenbreed}}

Werk je de ontwikkeling naar de eerste rij helemaal uit, dan krijg je een formule met zes producten van drie elementen. Die kun je onthouden met een schema.

\begin{theorie}
\begin{stelling}{Regel van Sarrus}
Voor een $3 \times 3$-matrix $A$ is

\vspace{1mm}
\qquad $\formule{\det A = a_{11}a_{22}a_{33} + a_{12}a_{23}a_{31} + a_{13}a_{21}a_{32}}$

\qquad\qquad\quad $\formule{{} - a_{13}a_{22}a_{31} - a_{11}a_{23}a_{32} - a_{12}a_{21}a_{33}}$
\end{stelling}

\begin{bewijs}
Ontwikkel naar de eerste rij:
\bewijsregel{\det A = a_{11}(a_{22}a_{33} - a_{23}a_{32}) - a_{12}(a_{21}a_{33} - a_{23}a_{31}) + a_{13}(a_{21}a_{32} - a_{22}a_{31}) .}
Werk de haakjes uit. De drie positieve termen zijn $a_{11}a_{22}a_{33}$, $a_{12}a_{23}a_{31}$ en $a_{13}a_{21}a_{32}$, de drie negatieve $a_{11}a_{23}a_{32}$, $a_{12}a_{21}a_{33}$ en $a_{13}a_{22}a_{31}$. Dat zijn precies de zes termen van de formule.
\end{bewijs}
\end{theorie}

Om de formule te onthouden, schrijf je de eerste twee kolommen nog eens rechts van de matrix. De drie positieve termen zijn de producten op de dalende diagonalen, de drie negatieve die op de stijgende diagonalen.

\naast{\begin{minipage}{9.4cm}
Voor de matrix $B$ uit het voorbeeld:
\[ \det B = (0 - 8 + 3) - (0 + 4 + 1) = -10 . \]
De dalende diagonalen geven $2 \cdot 0 \cdot (-1)$, $(-1) \cdot 2 \cdot 4$ en $3 \cdot 1 \cdot 1$, de stijgende $3 \cdot 0 \cdot 4$, $2 \cdot 2 \cdot 1$ en $(-1) \cdot 1 \cdot (-1)$. Sarrus werkt alleen bij $3 \times 3$-matrices: bij een $4 \times 4$-matrix geven de diagonalen niet de determinant.
\end{minipage}}{\begin{tikzpicture}[font=\small]
  \matrix[getallen,nodes={minimum width=6mm}] (m) {2 \& -1 \& 3 \& |[grijs]| 2 \& |[grijs]| -1 \\ 1 \& 0 \& 2 \& |[grijs]| 1 \& |[grijs]| 0 \\ 4 \& 1 \& -1 \& |[grijs]| 4 \& |[grijs]| 1 \\};
  \foreach \c in {1,2,3}{\pgfmathtruncatemacro{\e}{\c+2}\draw[accent,thick,opacity=.55] (m-1-\c.north west) -- (m-3-\e.south east);}
  \foreach \c in {3,4,5}{\pgfmathtruncatemacro{\b}{\c-2}\draw[fel!80!black,thick,dashed,opacity=.8] (m-1-\c.north east) -- (m-3-\b.south west);}
\end{tikzpicture}}

\begin{werkwijze}
Bij een $2 \times 2$-matrix: $ad - bc$. Bij een $3 \times 3$-matrix: ontwikkel naar een rij of kolom met veel nullen, met de tekens van het schaakbord, of gebruik de regel van Sarrus. Controleer met het rekentoestel als het mag.
\end{werkwijze}

\begin{oefeningen}
\oef{h3-det}
\end{oefeningen}

\subsection{De determinant van een product}

Hoe hangt de determinant van een product samen met die van de factoren? Voor $2 \times 2$-matrices kunnen we dat volledig uitrekenen.

\begin{theorie}
\begin{stelling}{Productregel}
Voor vierkante $n \times n$-matrices $A$ en $B$ is \qquad $\formule{\det(AB) = \det A \cdot \det B}$
\end{stelling}

\begin{bewijs}
We bewijzen het voor $n = 2$. Neem $A = \km{a & b \\ c & d}$ en $B = \km{p & q \\ r & s}$. Dan is
\bewijsregel{AB = \begin{bmatrix} ap + br & aq + bs \\ cp + dr & cq + ds \end{bmatrix}}
en
\bewijsregel{\det(AB) = (ap + br)(cq + ds) - (aq + bs)(cp + dr) .}
Bij het uitwerken vallen $apcq$ en $brds$ weg tegen $aqcp$ en $bsdr$. Er blijft
\bewijsregel{apds + brcq - aqdr - bscp = ad(ps - qr) - bc(ps - qr) = (ad - bc)(ps - qr) ,}
en dat is $\det A \cdot \det B$.
\end{bewijs}
\end{theorie}

Voor grotere matrices geldt de regel ook, maar dat bewijs is te zwaar voor deze cursus. De determinant van een product is dus het product van de determinanten. Voor een som bestaat zo'n regel niet: met $A = B = I_2$ is $\det(A + B) = \det(2I_2) = 4$, maar $\det A + \det B = 2$.

\voorbeeld{}{Bereken $\det(A^{10})$ voor $A = \begin{bmatrix} 2 & 3 \\ 1 & 1 \end{bmatrix}$.\label{vb:h3-detprod}}{%
\uitleg{$A^{10}$ uitrekenen is veel werk. Maar $A^{10}$ is een product van tien factoren $A$, dus volgens de productregel is de determinant een product van tien keer $\det A$.}
\begin{stappen}
\det A = 2 \cdot 1 - 3 \cdot 1 = -1 & De determinant van één factor. \\[1ex]
\det(A^{10}) = (\det A)^{10} = (-1)^{10} = 1 & De productregel, negen keer toegepast.
\end{stappen}}

\begin{oefeningen}
\oef{h3-k-detprod}
\end{oefeningen}

\subsection{Rijbewerkingen en de determinant}

Rijherleiden verandert een matrix, maar op een heel voorspelbare manier. Voor elk van de drie elementaire rijbewerkingen zegt een eenvoudige regel wat er met de determinant gebeurt. We bewijzen ze voor $2 \times 2$- en $3 \times 3$-matrices, met de stelling over het ontwikkelen; voor grotere matrices gelden ze ook.

\begin{theorie}
\begin{stelling}{Rijen verwisselen}
Verwissel je twee rijen, dan verandert de determinant van teken.
\end{stelling}

\begin{bewijs}
Voor $2 \times 2$-matrices: $\begin{vmatrix} c & d \\ a & b \end{vmatrix} = cb - da = -(ad - bc)$. Bij een $3 \times 3$-matrix ontwikkel je naar de rij die niet verwisseld wordt. Haar elementen en tekens blijven dezelfde, en in elke minor zijn twee rijen verwisseld: elke minor verandert van teken, en dus ook de determinant.
\end{bewijs}
\end{theorie}

Daaruit volgt meteen: heeft een matrix twee gelijke rijen, dan is haar determinant 0. Verwissel je die twee rijen, dan verandert de matrix niet, maar de determinant wel van teken. Dus is $\det A = -\det A$, en dat kan alleen als $\det A = 0$.

\begin{theorie}
\begin{stelling}{Een rij vermenigvuldigen}
Vermenigvuldig je één rij met een getal $r$, dan wordt de determinant $r$ keer zo groot.
\end{stelling}

\begin{bewijs}
Ontwikkel naar die rij (bij een $2 \times 2$-matrix zijn de minoren gewone getallen). Elk element van de rij wordt $r$ keer zo groot, en de minoren veranderen niet, want daarin is die rij geschrapt. Elke term, en dus de som, wordt $r$ keer zo groot.
\end{bewijs}
\end{theorie}

Let op: $rA$ is niet één rij maal $r$, maar elke rij. Bij een $3 \times 3$-matrix is dus $\det(2A) = 2 \cdot 2 \cdot 2 \cdot \det A = 8 \det A$, en niet $2 \det A$.

\begin{theorie}
\begin{stelling}{Een veelvoud van een rij optellen}
Tel je bij een rij een veelvoud van een andere rij op, dan verandert de determinant niet.
\end{stelling}

\begin{bewijs}
Stel dat je bij rij $i$ het $r$-voud van rij $j$ optelt. Ontwikkel naar rij $i$. De minoren veranderen niet, en element $k$ van die rij wordt $a_{ik} + ra_{jk}$. Splits elke term in twee: zo wordt de determinant
\bewijsregel{\det A + r \cdot D ,}
met $D$ de determinant van de matrix waarin rij $i$ vervangen is door rij $j$. Die matrix heeft twee gelijke rijen, dus $D = 0$.
\end{bewijs}
\end{theorie}

Die laatste regel is de nuttigste: met rijbewerkingen van de derde soort maak je nullen in een kolom, zonder dat de determinant verandert. Daarna ontwikkel je naar die kolom.

\voorbeeld{}{Bereken $\det A$ voor $A = \begin{bmatrix} 1 & 2 & -1 \\ 2 & 5 & 1 \\ -3 & -4 & 6 \end{bmatrix}$ met rijbewerkingen.\label{vb:h3-detrij}}{%
\uitleg{Er staat geen enkele 0 in de matrix. Met de 1 linksboven maken we eerst nullen in de eerste kolom: dat verandert de determinant niet. Dan blijft er bij het ontwikkelen maar één term over.}
\begin{stappenbreed}
\det A = \begin{vmatrix} 1 & 2 & -1 \\ 0 & 1 & 3 \\ 0 & 2 & 3 \end{vmatrix} \ \begin{array}{l} \\ \rb{R_2 - 2R_1} \\ \rb{R_3 + 3R_1} \end{array} & Een veelvoud van $R_1$ optellen: de determinant blijft gelijk. \\[5ex]
= 1 \cdot \begin{vmatrix} 1 & 3 \\ 2 & 3 \end{vmatrix} = 3 - 6 = -3 & Ontwikkel naar de eerste kolom: alleen de 1 levert een term.
\end{stappenbreed}
Controle met Sarrus: $(30 - 6 + 8) - (15 - 4 + 24) = 32 - 35 = -3$.}

\begin{werkwijze}
Maak met rijbewerkingen van de derde soort nullen in een kolom, liefst met een 1 als spil. Verwissel je rijen, verander dan het teken; deel je een rij door $r$, vermenigvuldig de determinant dan met $r$. Ontwikkel op het einde naar die kolom.
\end{werkwijze}

\begin{oefeningen}
\oef{h3-k-detrij}
\end{oefeningen}

\subsection{De determinant van de getransponeerde}

Zijn de rijen ook voor de determinant belangrijker dan de kolommen? Nee: transponeren verandert de determinant niet.

\begin{theorie}
\begin{stelling}{Getransponeerde}
Voor een vierkante matrix $A$ is \qquad $\formule{\det A\tr = \det A}$
\end{stelling}

\begin{bewijs}
Voor $2 \times 2$-matrices: $\begin{vmatrix} a & c \\ b & d \end{vmatrix} = ad - cb = \begin{vmatrix} a & b \\ c & d \end{vmatrix}$. Voor $3 \times 3$-matrices vervang je in de formule van Sarrus elke $a_{ij}$ door $a_{ji}$. De positieve termen worden $a_{11}a_{22}a_{33}$, $a_{21}a_{32}a_{13}$ en $a_{31}a_{12}a_{23}$, de negatieve $a_{31}a_{22}a_{13}$, $a_{11}a_{32}a_{23}$ en $a_{21}a_{12}a_{33}$: dezelfde zes termen als voor $A$.
\end{bewijs}
\end{theorie}

De kolommen van $A$ zijn de rijen van $A\tr$. Alles wat we over rijen zeiden, geldt dus ook voor kolommen: twee kolommen verwisselen verandert het teken, en een veelvoud van een kolom bij een andere optellen verandert niets.

\voorbeeld{}{Toon zonder veel rekenwerk aan dat $\begin{vmatrix} 1 & 4 & 2 \\ 2 & -1 & 4 \\ 3 & 5 & 6 \end{vmatrix} = 0$.\label{vb:h3-kolommen}}{%
\uitleg{De rijen lijken niets met elkaar te maken te hebben, maar de derde kolom is het dubbel van de eerste. Die kunnen we wegwerken met een kolombewerking. Kolom $j$ noemen we $K_j$.}
\begin{stappenbreed}
K_3 = 2K_1 & De derde kolom $2,\ 4,\ 6$ is twee keer de eerste $1,\ 2,\ 3$. \\[1ex]
\begin{vmatrix} 1 & 4 & 2 \\ 2 & -1 & 4 \\ 3 & 5 & 6 \end{vmatrix} = \begin{vmatrix} 1 & 4 & 0 \\ 2 & -1 & 0 \\ 3 & 5 & 0 \end{vmatrix} \ \rb{K_3 - 2K_1} & Een veelvoud van een kolom aftrekken verandert de determinant niet, want het is een rijbewerking op $A\tr$. \\[5ex]
= 0 & Ontwikkel naar de nulkolom: elke term is 0.
\end{stappenbreed}}

\subsection{Determinant en inverse}

Nu kunnen we bewijzen wat de inleiding beloofde: de determinant zegt of een matrix inverteerbaar is.

\begin{theorie}
\begin{stelling}{Inverteerbaar}
Een vierkante matrix $A$ is inverteerbaar als en slechts als $\det A \neq 0$. Dan is

\vspace{1mm}
\qquad $\formule{\det(A^{-1}) = \dfrac{1}{\det A}}$
\end{stelling}

\begin{bewijs}
Stel dat $A$ inverteerbaar is. Uit $AA^{-1} = I$ en de productregel volgt
\bewijsregel{\det A \cdot \det(A^{-1}) = \det I = 1 ,}
dus is $\det A \neq 0$ en $\det(A^{-1}) = \frac{1}{\det A}$.
Stel omgekeerd dat $\det A \neq 0$, en rijherleid $A$ tot een matrix $R$. Elke rijbewerking vermenigvuldigt de determinant met $-1$, met $r \neq 0$ of met 1, dus ook $\det R \neq 0$. Dan heeft $R$ geen nulrij, want ontwikkelen naar een nulrij geeft 0. Een vierkante rijgereduceerde matrix zonder nulrij heeft in elke rij een spil, en dus ook in elke kolom: ze is $I$. Volgens de stelling over de inverse met Gauss-Jordan is $A$ dan inverteerbaar.
\end{bewijs}
\end{theorie}

Bij $2 \times 2$-matrices geeft de formule meteen $A^{-1}$; bij grotere steunt het bewijs op eigenschappen die we niet algemeen bewezen. Een matrix met twee evenredige rijen heeft dus geen inverse.

\voorbeeld{}{Voor welke reële $k$ is $A = \begin{bmatrix} 1 & 0 & k \\ 0 & 1 & 2 \\ k & 2 & 5 \end{bmatrix}$ inverteerbaar?\label{vb:h3-invpar}}{%
\uitleg{We berekenen de determinant als uitdrukking in $k$, en zoeken waar ze 0 is. De eerste rij heeft een 0, dus ontwikkelen we daarnaar.}
\begin{stappenbreed}
\det A = 1 \cdot \begin{vmatrix} 1 & 2 \\ 2 & 5 \end{vmatrix} - 0 + k \cdot \begin{vmatrix} 0 & 1 \\ k & 2 \end{vmatrix} & Tekens $+$, $-$, $+$ in de eerste rij. \\[4ex]
= 1 \cdot (5 - 4) + k \cdot (0 - k) = 1 - k^2 & Twee $2 \times 2$-determinanten. \\[1ex]
1 - k^2 = 0 \iff k = 1 \text{ of } k = -1 & Daar is $A$ niet inverteerbaar.
\end{stappenbreed}
$A$ is dus inverteerbaar voor elke $k$ behalve $1$ en $-1$.}

\begin{oefeningen}
\oef{h3-k-invpar}

\oef{h3-eigensch}
\end{oefeningen}

\subsection{De regel van Cramer}

Bij een stelsel met evenveel vergelijkingen als onbekenden en een determinant verschillend van 0, kun je de oplossing ook met determinanten opschrijven.

\begin{theorie}
\begin{stelling}{Regel van Cramer}
Is $\det A \neq 0$, dan heeft $AX = B$ precies één oplossing, met \quad $\formule{x_j = \dfrac{\det A_j}{\det A}}$

Hierin is $A_j$ de matrix $A$ waarin kolom $j$ vervangen is door $B$.
\end{stelling}

\begin{bewijs}
We bewijzen het voor twee onbekenden; voor meer onbekenden is het bewijs te zwaar voor deze cursus. Neem $ax + by = e$ en $cx + dy = f$. Omdat $\det A = ad - bc \neq 0$, is $A$ inverteerbaar, en is de enige oplossing
\bewijsregel{\begin{bmatrix} x \\ y \end{bmatrix} = A^{-1}B = \frac{1}{ad - bc}\begin{bmatrix} d & -b \\ -c & a \end{bmatrix}\begin{bmatrix} e \\ f \end{bmatrix} = \frac{1}{ad - bc}\begin{bmatrix} de - bf \\ af - ce \end{bmatrix} .}
Daarin is $de - bf = \begin{vmatrix} e & b \\ f & d \end{vmatrix} = \det A_1$ en $af - ce = \begin{vmatrix} a & e \\ c & f \end{vmatrix} = \det A_2$.
\end{bewijs}
\end{theorie}

\voorbeeld{}{Los op met de regel van Cramer: $\begin{cases} 3x - 2y = 4 \\ 5x + y = 11 \end{cases}$\label{vb:h3-cramer}}{%
\uitleg{Eerst $\det A$, dan de determinanten met telkens één kolom vervangen door de rechterleden.}
\begin{stappen}
\det A = \begin{vmatrix} 3 & -2 \\ 5 & 1 \end{vmatrix} = 3 + 10 = 13 & Niet 0: precies één oplossing. \\[2ex]
x = \frac{1}{13}\begin{vmatrix} 4 & -2 \\ 11 & 1 \end{vmatrix} = \frac{4 + 22}{13} = 2 & Kolom 1 vervangen door $\km{4 \\ 11}$. \\[2ex]
y = \frac{1}{13}\begin{vmatrix} 3 & 4 \\ 5 & 11 \end{vmatrix} = \frac{33 - 20}{13} = 1 & Kolom 2 vervangen.
\end{stappen}}

Is $\det A = 0$, dan zegt de regel niets: $x + y = 1$, $2x + 2y = 2$ is onbepaald en $x + y = 1$, $2x + 2y = 3$ strijdig. Bij grote stelsels is Gauss-Jordan minder werk.

\begin{oefeningen}
\oef{h3-cramer}
\end{oefeningen}

%========================================================================================
\section{De rang}

\subsection{Definitie}

In een stelsel kan een vergelijking een gevolg zijn van de andere, zoals in voorbeeld \ref{vb:h2-onbepaald} van hoofdstuk 2: na rijherleiden werd ze een nulrij. De rang telt hoeveel vergelijkingen er overblijven.

\begin{theorie}
De \textbf{\textcolor{accent}{rang}} van een matrix $A$ is het aantal spillen in de rijgereduceerde vorm van $A$. Je noteert $\formule{\operatorname{rang} A}$.
\end{theorie}

Het aantal spillen is ook het aantal rijen dat niet nul is in de rijgereduceerde vorm. Omdat elke spil in een eigen rij en een eigen kolom staat, is de rang van een $m \times n$-matrix hoogstens het kleinste van $m$ en $n$.

\voorbeeld{}{Bepaal de rang van $R = \begin{bmatrix} 1 & 2 & -1 & 3 \\ 2 & 4 & 0 & 2 \\ 3 & 6 & -1 & 5 \end{bmatrix}$.\label{vb:h3-rang}}{%
\uitleg{We rijherleiden en tellen de spillen.}
\begin{stappenbreed}
\begin{bmatrix} 1 & 2 & -1 & 3 \\ 0 & 0 & 2 & -4 \\ 0 & 0 & 2 & -4 \end{bmatrix} \ \begin{array}{l} \\ \rb{R_2 - 2R_1} \\ \rb{R_3 - 3R_1} \end{array} & Kolom 2 krijgt geen spil: onder de 2 staan alleen nullen. \\[5ex]
\begin{bmatrix} 1 & 2 & 0 & 1 \\ 0 & 0 & 1 & -2 \\ 0 & 0 & 0 & 0 \end{bmatrix} \ \begin{array}{l} \rb{R_1 + \frac12 R_2} \\ \rb{\frac12 R_2} \\ \rb{R_3 - R_2} \end{array} & Twee spillen: $\operatorname{rang} R = 2$.
\end{stappenbreed}
De derde rij van $R$ is de som van de eerste twee: ze brengt niets nieuws.}

\begin{oefeningen}
\oef{h3-rang}
\end{oefeningen}

\subsection{Rang en stelsels}

Met de rang vat je samen wat je in hoofdstuk 2 over de soorten stelsels leerde: je vergelijkt de rang van de coëfficiëntenmatrix met die van de uitgebreide matrix.

\begin{theorie}
\begin{stelling}{Rang en stelsels}
Een stelsel $AX = B$ met $n$ onbekenden is strijdig als en slechts als $\formule{\operatorname{rang} A < \operatorname{rang}\,[A \,|\, B]}$. Is het niet strijdig, dan heeft het $\formule{n - \operatorname{rang} A}$ parameters, en is het bepaald als dat 0 is.
\end{stelling}

\begin{bewijs}
Rijherleid $[A \,|\, B]$. Het deel links van de streep is dan de rijgereduceerde vorm van $A$, want dezelfde rijbewerkingen herleiden $A$. Links staan dus $\operatorname{rang} A$ spillen, en $[A \,|\, B]$ heeft er hoogstens één meer: een spil in de laatste kolom. Zo'n spil is een rij $[0 \ \cdots \ 0 \,|\, 1]$, en volgens hoofdstuk 2 is het stelsel strijdig precies als zo'n rij er is. Is er geen, dan heeft elke onbekende zonder spil een parameter, en dat zijn er $n - \operatorname{rang} A$.
\end{bewijs}
\end{theorie}

\voorbeeld{}{Bepaal $\operatorname{rang} A$ en $\operatorname{rang}\,[A \,|\, B]$ voor het stelsel uit voorbeeld \ref{vb:h2-strijdig} van hoofdstuk 2. Wat besluit je?\label{vb:h3-rangstelsel}}{%
\uitleg{De rijherleiding staat al in hoofdstuk 2. We tellen de spillen links van de streep en in de hele matrix.}
\begin{stappenbreed}
\um{ccc|c}{1 & 0 & -7 & -1 \\ 0 & 1 & 3 & 2 \\ 0 & 0 & 0 & 1} & Na de rijbewerkingen uit hoofdstuk 2. Nullen boven de laatste 1 veranderen het aantal spillen niet meer. \\[5ex]
\operatorname{rang} A = 2, \quad \operatorname{rang}\,[A \,|\, B] = 3 & Links twee spillen, en in de laatste kolom nog een derde. \\[1ex]
\operatorname{rang} A < \operatorname{rang}\,[A \,|\, B] & Het stelsel is strijdig.
\end{stappenbreed}}

\nujij{Een stelsel met 4 onbekenden heeft een coëfficiëntenmatrix met rang 3, en de uitgebreide matrix heeft ook rang 3. Hoeveel oplossingen zijn er?}{Niet strijdig, en $4 - 3 = 1$ parameter: oneindig veel oplossingen.}

\begin{oefeningen}
\oef{h3-rangst}
\end{oefeningen}

\subsection{Vierkante matrices}

Bij een vierkante matrix komen rang, inverse, determinant en stelsels samen in één stelling.

\begin{theorie}
\begin{stelling}{Rang, determinant en inverse}
Voor een vierkante $n \times n$-matrix $A$ zijn gelijkwaardig:

\vspace{1mm}
\qquad $\formule{\operatorname{rang} A = n}$ \ $\iff$ \ $A$ is inverteerbaar \ $\iff$ \ $\formule{\det A \neq 0}$ \ $\iff$ \ $AX = B$ is bepaald voor elke $B$.
\end{stelling}

\begin{bewijs}
Een $n \times n$-matrix heeft rang $n$ als en slechts als haar rijgereduceerde vorm in elke rij en elke kolom een spil heeft, dus als ze rijherleidt tot $I$. Volgens de stelling over de inverse met Gauss-Jordan is dat zo als en slechts als $A$ inverteerbaar is, en dat is zo als en slechts als $\det A \neq 0$.
Is $\operatorname{rang} A = n$, dan is $\operatorname{rang}\,[A \,|\, B] = n$ voor elke $B$, want er zijn maar $n$ rijen. Het stelsel is dan niet strijdig en heeft $n - n = 0$ parameters: het is bepaald. Is $\operatorname{rang} A < n$, neem dan $B = O$. Het stelsel $AX = O$ heeft de oplossing $X = O$, is dus niet strijdig, en heeft $n - \operatorname{rang} A > 0$ parameters: het is niet bepaald.
\end{bewijs}
\end{theorie}

Een vierkante matrix is dus ofwel \textit{goed}: rang $n$, inverteerbaar, determinant niet 0, en elk stelsel heeft precies één oplossing. Ofwel is ze dat allemaal niet. Om te weten in welk geval je zit, volstaat één van de vier: kies wat het minste rekenwerk vraagt.

\voorbeeld{}{Is $A = \begin{bmatrix} 1 & 2 & 3 \\ 4 & 5 & 6 \\ 7 & 8 & 9 \end{bmatrix}$ inverteerbaar? Wat is de rang van $A$, en hoeveel oplossingen heeft $AX = O$?\label{vb:h3-vierkant}}{%
\uitleg{De determinant is het snelst. Is ze 0, dan weten we meteen dat de rang kleiner is dan 3.}
\begin{stappenbreed}
\det A = (45 + 84 + 96) - (105 + 48 + 72) = 0 & Sarrus. Dus $A$ is niet inverteerbaar. \\[1ex]
\begin{bmatrix} 1 & 2 & 3 \\ 0 & -3 & -6 \\ 0 & -6 & -12 \end{bmatrix} \ \begin{array}{l} \\ \rb{R_2 - 4R_1} \\ \rb{R_3 - 7R_1} \end{array} & De laatste rij is het dubbel van de middelste: na $R_3 - 2R_2$ een nulrij. $\operatorname{rang} A = 2$. \\[5ex]
3 - 2 = 1 & $AX = O$ heeft één parameter: oneindig veel oplossingen, en niet alleen $X = O$.
\end{stappenbreed}}

\begin{oefeningen}
\oef{h3-k-vierkant}
\end{oefeningen}

%========================================================================================
\opdrachtenkop

\begin{oefeningen}
\needspace{14\baselineskip}
\oef{h3-o14}

\oef{h3-o17}

\oef{h3-o4}

\needspace{14\baselineskip}
\oef{h3-o15}

\oef{h3-o19}

\oef{h3-o2}

\oef{h3-o16}

\needspace{14\baselineskip}
\oef{h3-o18}

\oef{h3-o1}

\oef{h3-o3}

\needspace{14\baselineskip}
\oef{h3-o21}

\oef{h3-o5}

\needspace{14\baselineskip}
\oef{h3-o24}

\oef{h3-o7}

\oef{h3-o8}

\needspace{14\baselineskip}
\oef{h3-o20}

\oef{h3-o22}

\oef{h3-o6}

\oef{h3-o9}

\oef{h3-o23}

\oef{h3-o10}

\oef{h3-o25}

\oef{h3-o12}

\oef{h3-o13}

\oef{h3-o26}

\oef{h3-o11}
\end{oefeningen}

%----------------------------------------------------------------------------------------
% KOPPEL: verbanden met andere leerstof
%----------------------------------------------------------------------------------------
\needspace{20\baselineskip}
\opdrachtsectie{Koppel}

\lkr{Coördinaten en de oppervlakte van een driehoek kennen de leerlingen uit de 2de graad. Complexe getallen staat in hetzelfde jaar als deze cursus: kijk na of die cursus al gezien is.}

\begin{oefeningen}
\oef{h3-koppel-1}

\oef{h3-koppel-2}
\end{oefeningen}

%----------------------------------------------------------------------------------------
\opdrachtsectie{Test jezelf}

\begin{oefeningen}
\oef{h3-t1}

\oef{h3-t2}

\oef{h3-t3}

\oef{h3-t4}

\oef{h3-t5}

\oef{h3-t6}

\oef{h3-t7}

\oef{h3-t8}
\end{oefeningen}

\herhaalschema{%
  \fv{oef:h3-t1}{vb:h3-inverse2}{oef:h3-inverse}%
  \fv{oef:h3-t2}{vb:h3-inverse}{oef:h3-inverse}%
  \fv{oef:h3-t3}{vb:h3-det}{oef:h3-det}%
  \fv{oef:h3-t4}{vb:h3-cramer}{oef:h3-cramer}%
  \fv{oef:h3-t5}{vb:h3-rang}{oef:h3-rang}%
  \fv{oef:h3-t6}{vb:h3-invpar}{oef:h3-k-invpar}%
  \fv{oef:h3-t7}{vb:h3-vgl}{oef:h3-matrixvgl}%
  \fv{oef:h3-t8}{vb:h3-rangstelsel}{oef:h3-rangst}}

\ifzelfstudie
\section*{Maak je eigen samenvatting}
Sluit dit boek en schrijf in je schrift: wat een inverse matrix is en waarom ze uniek is, hoe je een matrixvergelijking oplost, hoe je de inverse berekent met $[A \,|\, I]$, de determinant en de inverse van een $2 \times 2$-matrix, hoe je een determinant van een $3 \times 3$-matrix berekent, de eigenschappen van de determinant, de regel van Cramer, wat de rang is, en het verband tussen rang, determinant, inverse en stelsels. Vergelijk daarna met de kaders.
\fi
